\subsection*{Derivation of the precision optimum}
\label{sec:analytical:deriv}

We derive analytically the precision level $\pstar$ that maximizes metacognitive accuracy
(Equation~\ref{eq:meta}). Substituting $\uncertainty = 1 - \precision$:

\begin{align}
  \metaacc(\precision) &= \frac{1}{2}\left[
    \sigmoid\!\bigl((1-\precision)\slope - \thresh{\text{h}}\bigr) +
    1 - \sigmoid\!\bigl((1-\precision)\slope - \thresh{\text{e}}\bigr)
  \right]
\end{align}

Consider the behavior at the extremes. As $\precision \to 1$ (maximal precision, minimal uncertainty):
\begin{align}
  (1-\precision)\slope \to 0, \quad
  \sigmoid(0 - \thresh{\text{h}}) &= \sigmoid(-\thresh{\text{h}}) \approx 0.27, \quad
  \sigmoid(0 - \thresh{\text{e}}) = \sigmoid(-\thresh{\text{e}}) \approx 0.076
\end{align}
So $\metaacc \to (0.27 + 0.924)/2 = 0.597$. Intuitively, an overconfident agent
never feels uncertain enough to opt out on hard trials, so $\poptout(\text{hard}) \to
\sigmoid(-\thresh{\text{h}}) < 0.5$.

As $\precision \to 0$ (minimal precision, maximal uncertainty):
\begin{align}
  (1-\precision)\slope \to \slope, \quad
  \sigmoid(\slope - \thresh{\text{h}}) &\to 1, \quad
  \sigmoid(\slope - \thresh{\text{e}}) \to 1
\end{align}
So $\metaacc \to (1 + 0)/2 = 0.5$. An agent with no precision opts out on
\emph{every} trial, including easy ones, achieving only chance-level accuracy.

By continuity and the intermediate value theorem, a maximum $\pstar \in (0,1)$ must exist.
Numerically, with $\slope = 4.0$, $\thresh{\text{h}} = 1.0$, $\thresh{\text{e}} = 2.5$:

\begin{equation}
  \boxed{\pstar = 0.563, \quad \metaacc(\pstar) = 0.679}
\end{equation}

Figure~\ref{fig:inverted_u} shows the full precision--metacognition curve and the sensitivity of
$\pstar$ to slope variation.

\subsection*{Robustness}
\label{sec:analytical:robust}

The qualitative result---an inverted-U with an interior maximum---holds for any slope $\slope \in [2, 6]$
and any threshold separation $\thresh{\text{e}} - \thresh{\text{h}} > 0$. The exact value of $\pstar$
shifts modestly with slope (range $\pstar \in [0.50, 0.63]$), but the directional prediction remains:
\emph{intermediate precision is optimal for metacognitive accuracy}. The robustness band is shown as the
shaded region in Figure~\ref{fig:inverted_u}.

\subsection*{Interpretation: the overconfidence paradox}
\label{sec:analytical:paradox}

The overconfidence paradox can be stated precisely: increasing precision improves performance on
Domains 2--4 monotonically (through the $\precision \cdot \text{weight}$ structure), but
\emph{degrades} metacognitive accuracy for $\precision > \pstar$. This creates an unavoidable
architectural tension: the same precision parameter that drives cognitive excellence undermines
self-knowledge. In ecological terms, a highly accurate forager may be less able to
``know what it doesn't know.''

This result has a formal connection to signal detection theory: an agent that places very tight priors
on sensory signals (high precision) will classify ambiguous stimuli as certain, reducing the
uncertainty signal available for opt-out decisions. The inverted-U is thus not a pathology of the
model but a fundamental property of precision-weighted inference under threshold-based opt-out
policies.
